% =============================================
% BYLAWS CONTENT FILE
% Edit ONLY this file for content changes.
% Use normal LaTeX enumerate environments for all numbered/lettered lists.
% The formatting is handled automatically in main.tex
% =============================================

\JBArticle {pb-article-1} {Article I: Purpose and Authority}
\begin {enumerate}
    \item The Parking Board is a subordinate judicial body of the Judicial Board, established as a partnership between the Rensselaer Union and the Director of Auxiliary Services (DAS), to oversee and adjudicate appeals related to parking citations.
    \item A contestation of a parking ticket does not constitute an appeal; however, all parking citations maintain the right to come before this Board for a formal appeal.
    \item The Parking Board explicitly recognizes the superior authority of the Institute Hearing Procedures.
    \item Judgments rendered by the Parking Board must follow Institute Hearing Procedures and shall be adhered to by the administration. 
    \item Appeals of this Board's decisions may be sought only if they demonstrate grounds for appeal as defined in the RPI Student Handbook of Rights and Responsibilities.
\end {enumerate}

\JBArticle {pb-article-2} {Article II: Membership and Composition}
\begin {enumerate}
    \item The Parking Board is an ad hoc committee that remains omnipresent in its authority but is fully formed only when notified of a pending case by the DAS or their designee.
    \item A standard panel consists of four (4) voting members:
    \begin {enumerate}
        \item Two (2) members or alternates of the Judicial Board, selected by the Parking Board Chairperson.
        \item Two (2) administrators appointed by the DAS or their designee.
    \end {enumerate}
    \item The Chairperson shall only cast a vote in the event of a tie between the other four members.
\end {enumerate}

\JBArticle {pb-article-3} {Article III: Quorum and Administrative Absence}
\begin {enumerate}
    \item In the event that the DAS fails to appoint administrative representatives, the responsibility to ensure a full panel falls to the Parking Board Chairperson.
    \item The Chairperson may appoint any employee of the Institute who is not a student to serve as an administrative representative, provided the following conditions are met:
    \begin {enumerate}
        \item The Chairperson of the Judicial Board provides written approval of the appointment.
        \item The appointee provides written consent expressing their desire to serve.
    \end {enumerate}
\end {enumerate}

\JBArticle {pb-article-4} {Article IV: Burden of Proof and Standard of Evidence}
\begin {enumerate}
    \item In accordance with Institute Hearing Procedures, the standard of proof for all proceedings is a preponderance of the evidence (51\%).
    \item The burden of proof rests with the student appellant to demonstrate that the citation was not justly issued.
    \item For the purposes of this Board, "justly" shall be interpreted to mean the citation was both technically accurate and right to be issued, taking into account the specific circumstances of the incident.
\end {enumerate}

\JBArticle {pb-article-5} {Article V: Conflicts of Interest}
\begin {enumerate}
    \item The Chairperson maintains sole authority to determine if a conflict of interest exists.
    \item If a conflict is identified, the Chairperson may reselect student members or refuse administrative appointees.
    \item If an administrative appointee is conflicted, the DAS must provide a new appointee or serve on the panel themselves.
\end {enumerate}

\JBArticle {pb-article-6} {Article VI: Records and Confidentiality}
\begin {enumerate}
    \item Case records for parking citations are sealed by default. Only generalized statistics may be released to individuals outside the Board.
    \item Full, unamended case details shall be accessible only to the membership of the Judicial Board and the staff of the Institute Department of Auxiliary Services.
\end {enumerate}